
This pilot study presents the first measurement of logit-based Expected Calibration Error on adversarial NLP benchmark splits for an open-weight LLM, filling a gap between adversarial robustness evaluation (accuracy-only) and calibration measurement (clean-data-only). The central finding is that adversarial calibration degradation in LLMs is construction-method and task-type conditional, not universal.

When adversarial examples are constructed model-in-the-loop---selected precisely because the model fails on them---the resulting benchmark tends to contain examples where the model is already appropriately uncertain at failure time. Calibration improvement under these conditions follows directly. Calibration degrades in the case where adversarial examples are constructed human-adversarially, targeting surface features to trigger high model confidence on incorrect answers, independent of model uncertainty. AdvGLUE MNLI is this case: $\Delta ECE$ = +0.071, a 7.1 percentage-point calibration gap that renders confidence signals less reliable under human-crafted adversarial NLI conditions.

RLHF alignment interacts with the construction-method dependency: RLHF-aligned models show consistently better calibration than base models on ANLI ($\Delta \Delta ECE$ up to +0.147 across all 3 rounds) but worse calibration on AdvGLUE MNLI (alignment tax: $\Delta \Delta ECE = -0.026)$. RLHF calibration properties measured on clean benchmarks do not transfer to adversarial settings in a construction-method-independent way.

\textbf{Summary of main contributions:}

\begin{enumerate}
\item \textbf{First adversarial NLP calibration measurement.} $\Delta ECE$ = +0.071 for AdvGLUE MNLI confirms that adversarial calibration degradation is real and measurable. Label preservation rate = 1.000 for all adversarial splits validates $\Delta ECE$ as a genuine calibration signal.
\end{enumerate}

\begin{enumerate}
\item \textbf{Conditional structure: construction method and task type determine calibration outcome.} Human-adversarial NLI ($\Delta ECE$ = +0.071) degrades calibration; model-in-loop NLI (ANLI R1/R2: $\Delta ECE$ < 0) improves it. The ANLI difficulty gradient (R1 $\rightarrow$ R2 $\rightarrow$ R3: $-0.041 \rightarrow -0.014 \rightarrow$ +0.024) shows a smooth transition. Binary classification consistently shows reversed $\Delta ECE$.
\end{enumerate}

\begin{enumerate}
\item \textbf{Benchmark-type $\times$ RLHF interaction.} RLHF alignment conditionally moderates adversarial calibration degradation: confirmed for model-in-loop adversarial (ANLI: 100\% moderation rate, 3/3 rounds); reversed for static human-adversarial benchmarks (AdvGLUE: $\Delta \Delta ECE = -0.026)$.
\end{enumerate}

\textbf{Future directions.} The adaptive uncertainty hypothesis for ANLI R1/R2 calibration improvement can be directly tested by computing ECE separately for correct and incorrect predictions. The alignment tax mechanism can be tested by constructing human-adversarial NLI examples targeting RLHF models specifically. Completing the planned multi-model grid (Mistral-7B-Instruct, Llama-2-13b-chat) would enable proper evaluation of calibration universality and deployment reliability prediction. Temperature scaling applied post-hoc to existing logit caches would clarify whether calibration degradation under human-adversarial conditions is a confidence scaling artifact or a structural property.

Adversarial accuracy evaluation and calibration evaluation require shared vocabulary and measurement methods. The conditional structure documented here---where the same adversarial benchmark can improve or degrade calibration depending on construction method and RLHF alignment---suggests that these communities must develop coordinated frameworks. Calibration robustness is not a consequence of accuracy robustness; it requires independent measurement.

